% \cleardoublepage
\appendix
\section{Appendices}



\subsection{Data Augmentation}
During the data augmentation process, we mainly utilized two of the most popular visualization libraries: D3 \cite{d3} and Vegalite \cite{vegalite}. Moreover, we have introduced a range of visual variations in terms of color scheme, elements dimensions, shapes, background, .etc (see \Cref{fig:d3-samples}).  This makes our generated chart images closely resemble the real-world charts found on the Web. 

\begin{table*}
 \setlength\extrarowheight{2pt}
 \centering
 \caption{Chart type distribution and linguistic statistics of the chart pertaining corpus. The charts in the last group ({\color{magenta} magenta}) do not come with an underlying data table. The charts generated by the data augmentation process are shown in {\textcolor{azure}{blue}}. 
 }
 
 \scalebox{0.75}{\begin{tabular}{lcccccc?ccc}
  
  \toprule
  
   \multirow{3}{*}{} & \multicolumn{6}{c}{Chart Type} & \multicolumn{3}{c}{Linguistic Statistics}  \\ \cmidrule(lr){2-7} \cmidrule(lr){8-10} 
  
   Datasets & \multicolumn{2}{c}{Bar} & \multicolumn{2}{c}{Line} & \multirow{2}{*}{Pie} & \multirow{2}{*}{\#Charts} & \multirow{2}{*}{\#Vocab } & \multirow{2}{*}{Avg. Character} & \multirow{2}{*}{Avg. Token}\\ \cmidrule(lr){2-3} \cmidrule(lr){4-5}
                                             
    & Two-Col & Multi-Col & Two-Col & Multi-Col & & & &  &  \\ 
   \midrule
   % \cline{7-9}
  
  % \hline
  
  % Statista
  Statista & 71.9\% & 15.3\% & 8.3\% & 2\% & 2\% & \textbf{19,143} & 24,392 & 111.37 & 21.88\\ 
 
  % OWID
  OWID & 51.9\% & 0.0\% & 9\% & 38.9\% & 0.0\% & \textbf{60,624} & 3,721 & 85.89 & 16.96\\ 

  % OECD
  OECD & 49.1\% & 0.0\% & 3.1\%& 47.7\% & 0.0\% & \textbf{64,380} & 1,606 & 65.47 & 14.67\\ 

  % PlotQA
  PlotQA & 11.2\% & 55.6\% & 6.7\% & 26.2\% & 0.0\% & \textbf{157,070} & 2,230 & 155.32 & 33.21\\ 

  % Beagle
  Beagle & 29.8\% & 27.3\% & 24.7\% & 17.9\% & 0.0\% & \textbf{3,972} & 11,361 & 78.76 & 20.55\\ 

  % ChartInfo
  ChartInfo & 31.7\% & 51.0\% & 8.6\% & 8.6\% & 0.0\% & \textbf{1,796} & 13,329 & 120.75 & 26.11\\ 

  \midrule

% Data Augmentation
  
  \textcolor{azure}{Data Augmentation} & 13.3\% & 49.3\% & 11.7\% & 11.1\% & 14.3\% & \textbf{189,836} & 117,244 & 85.62 & 21.16 \\ 
  \midrule

  % ChartOCR
  {\color{magenta} ExcelChart400K} & 11.5\% & 32.7\% & 12.0\% & 22.3\% & 27.7\% & \textbf{106,897} & 515,922 & 138.68 & 27.72\\ 
  

  % PewResearch
  {\color{magenta} PewResearch} & 11.4\% & 55.5\% & 4.4\% & 21.9\% & 6.5\% & \textbf{5,295} & 38,165 & 477.33 & 98.08\\ 

  % LineCap
{\color{magenta} LineCap} & 0.0\% & 0.0\% & 15.9\% & 84.0\% & 0.0\% & \textbf{2,821} & 16,570 & 102.11 & 24.62\\ 

 % \rowcolor{gray!30} 
  % Neural Caption Generation
  {\color{magenta} Neural Caption} & 0.0\% & 0.0\% & 100\% & 0.0\% & 0.0\% & \textbf{100} & 389 & 117.56 & 27.43\\ 

  \midrule
  

  \textit{\textbf{Total}} & \textbf{21.95\%} & \textbf{36.56\%} & \textbf{9.23\%} & \textbf{23.71\%} & \textbf{9.39\%} & \textbf{611,934} & \textbf{888,522} & \textbf{114.70} & \textbf{25.01}\\ \bottomrule
  
  \hline
  
 \end{tabular}}
 \label{tab:datasets}
\end{table*}


\begin{table}[t!]
 \setlength\extrarowheight{2pt}
 \centering
 
 
 
 \scalebox{0.65}{\begin{tabular}{l|cccc}
  
  \toprule

   \textbf{Datasets} & \textbf{\#Vocab} & \textbf{Avg. Char} & \textbf{Avg. Token} & \textbf{Avg. Sentence}\\
   % \cmidrule(lr){1-5}
    \midrule                                        

  

  Statista & 72,725 & 450.28 & 106.68 & 4.46 \\  

  % OWID
  OWID & 58,212 & 463.48 & 105.99 & 4.47\\ 

  % OECD
  OECD & 24,752 & 414.97 & 95.08 & 4.78\\ 
  
  % PlotQA
  PlotQA & 112,394 & 666.09 & 149.84 & 5.20\\ 

  % DataAug
  Data Aug. & 162,239 & 468.41 & 113.46 & 4.49\\ 

  % PewResearch
  PewResearch & 13,449 & 604.04 & 133.01 & 4.66\\ 

  % LineCap
  LineCap & 2018 & 110.82 & 26.24 & 1.87\\ 

  % Neural Caption Generation
  Neural Caption & 1338 & 262.84 & 53.28 & 3.58\\ \bottomrule
  
  \hline
  
 \end{tabular}}
 \vspace{-3mm}
 \caption{\small{Statistics about the captions of the datasets used in LM pretraining. 
 }
 \vspace{-4mm}
 }
 
 \label{tab:textstats}
 
 \end{table}

 



\begin{figure*}[th]
\begin{subfigure}[b]{.49\textwidth}
\centering
    \includegraphics[width=0.8\textwidth,keepaspectratio]{Figures/WDCD3_Corpus10_33360.png}
\caption{}
\label{fig:d1}
\end{subfigure}
\begin{subfigure}[b]{.49\textwidth}
\centering
    \includegraphics[width=0.8\textwidth,keepaspectratio]{Figures/WDCD3_Corpus10_30959.png}
\caption{}
\label{fig:d2}
\end{subfigure}
\begin{subfigure}[b]{.49\textwidth}
\centering
    \includegraphics[width=0.8\textwidth,keepaspectratio]{Figures/WDCD3_Corpus10_33389.png}
\caption{}
\label{fig:d3}
\end{subfigure}
\begin{subfigure}[b]{.49\textwidth}
\centering
    \includegraphics[width=0.8\textwidth,keepaspectratio]{Figures/WDCD3_Corpus10_33648.png}
\caption{}
\label{fig:d4}
\end{subfigure}
\begin{subfigure}[b]{.49\textwidth}
\centering
    \includegraphics[width=0.8\textwidth,keepaspectratio]{Figures/WDCVegalite_corp11_917.png}
\caption{}
\label{fig:d5}
\end{subfigure}
\begin{subfigure}[b]{.49\textwidth}
\centering
    \includegraphics[width=0.8\textwidth,keepaspectratio]{Figures/WDCVegalite_corp11_946.png}
\caption{}
\label{fig:d6}
\end{subfigure}

\begin{subfigure}[b]{.49\textwidth}
\centering
    \includegraphics[width=0.55\textwidth,keepaspectratio]{Figures/WDCD3_Corpus10_17396.png}
\caption{}
\label{fig:d9}
\end{subfigure}
\begin{subfigure}[b]{.49\textwidth}
\centering
    \includegraphics[width=0.55\textwidth,keepaspectratio]{Figures/WDCD3_Corpus10_17431.png}
\caption{}
\label{fig:d10}
\end{subfigure}
\vspace{-1em}
\caption{Visually diverse charts generated by D3 and Vegalite for WDC corpus (\Cref{fig:d1}, \Cref{fig:d2}, \Cref{fig:d3}, \Cref{fig:d4}, \Cref{fig:d9} and \Cref{fig:d10} from D3-WDC, \Cref{fig:d5} and \Cref{fig:d6} from Vegalite-WDC.
% , \Cref{fig:d7} from D3-Nationmaster and \Cref{fig:d8} from Vegalite-Nationmaster). 
Visual factors like color scheme, width of bars, and existence of grids and axis labels are different among the samples.}
\label{fig:d3-samples}
\end{figure*}






\subsection{Data Augmentation by Knowledge Distillation}
\label{app:knowledge-distillation}

 We select the small dataset (3700 charts) from PlotQA and the augmented charts from WDC, since these datasets are accompanied by the underlying data tables which serve as suitable chart representation for LLMs. Also, they cover a wide range of topics which contributes to the diversity in the generated summaries (\Cref{subsec:charts-data}). \Cref{fig:instructgpt_summary} shows our process for generating summaries for the charts that have underlying data tables using InstructGPT model \cite{instructgpt}. The input mainly consists of one demonstration (table-caption pair) followed by the desired chart data table. The output is the generated summary. Using this mechanism, we generated a small dataset of 3,700 samples. We then finetuned Flan-T5 XL \cite{chung2022scaling} on this dataset. To our knowledge, Flan-T5 was the SoTA open-sourced instruction-tuned model during the development of our dataset. After finetuning on our task, we (qualitatively) observed similar performance as \texttt{text-davinci-003}. At the final step, we used the finetuned Flan-T5 model to generate summaries for all the charts that do not have an associated summary (e.g., PlotQA, augmented charts, OWID and OECD charts). In this process, we added around 470K  summaries for charts in our pretraining corpus. \Cref{fig:flant5_summaries} shows some examples generated by the finetuned Flan-T5.
 
 

To benefit more from the capability of GPT models in generating high-quality summaries, we further prompt ChatGPT (\texttt{gpt-3.5-turbo}) \cite{ChatGPT} to generate summaries for the charts from Statista and Pew Research and put these in our pretraining corpus instead of the original summaries in the Chart-to-Text benchmark \cite{chart-to-text-acl}. In most cases, we found the summaries from ChatGPT to be more elaborate and of better writing style.  For the Pew Research Centre charts, the underlying data tables are not provided. However, we have observed that the underlying data values are written on the visual elements in most of these charts. Hence, we decided to use an OCR tool to extract the layout-preserving texts from the chart images, and then feed it into ChatGPT to generate the summaries as shown in \Cref{fig:pew-layout-preserved}. We realized that ChatGPT is capable of understanding a chart from the OCR data.





 \begin{figure*}
    \centering
    \includegraphics[width=\textwidth]{Figures/InstructGPT_summary.png}
    \caption{An example of the performance of \textit{InstructGPT} in generation summaries for data tables. On the left side, the red text is a full example of a demonstration and its summary followed by the demonstration for the target chart. The paragraph in green shows the summary generated by the model.}
    \label{fig:instructgpt_summary}
\end{figure*}


\begin{figure*}
     \centering        \includegraphics[width=.7\textwidth]{Figures/inference_time_plot.png}
    \caption{Average inference time for 10 random samples from three major benchmarks in chart understanding domain for \model\ and MatCha models}
    \label{fig:inference_time}
\end{figure*} 


\begin{figure*}
    \centering
    \includegraphics[width=\textwidth]{Figures/flant5_summaries.drawio.pdf}
    \caption{Examples of summaries generated by Flan-T5 XL model after fine-tuning.}
    \label{fig:flant5_summaries}
\end{figure*}

 \begin{figure*}
    \centering
    \includegraphics[width=\textwidth]{Figures/Pew_ChatGPT.drawio1.pdf}
    \caption{An example of the layout-preserved OCR-extracted text for a PewResearch chart image where the underlying data table is not available. The extracted text is then given to ChatGPT to generate a summary. ChatGPT can still extract and comprehend important information and insights from the layout-preserving text of the chart image.}
    \label{fig:pew-layout-preserved}
\end{figure*}



\subsection{Dataset Analysis}
\label{app:dataset-analysis}
The linguistic characteristics of the textual elements vary across different datasets, with charts from PlotQA and PewResearch often having longer text elements (e.g., axis labels, legends, titles), while augmented data and Beagle datasets contain shorter text
 
 (\Cref{tab:datasets}, right). In \Cref{tab:textstats}, we further provide linguistic statistics for the summaries of the datasets used in the summary generation task at pretraining.


\subsection{Experiments Setup}
\label{app:details-downstream-tasks}
To minimize the computational resource requirements, we initialize our model from the base Donut weights \cite{kim2022ocr}. Our pretraining process consists of two stages. In the first stage, we set the input image resolution to 512x512 and pretrain for 300K steps. In the second stage, we increase the input image resolution to 960x960 and pretrain for an additional 100K steps. \Cref{tab:finetuning-details} shows the hyperparameters we used in pretraining and finetuning our model on each downstream task. All our experiments were carried out using one 4-A100 (40GB), one 4-A100 (80GB), and one 4-V100 (32 GB) GPU machines.  

\begin{table*}[t]
 \setlength\extrarowheight{2pt}
 \centering
 \caption{
  Training details for pretraining and finetuning experiments.
 }
 \scalebox{0.85}{\begin{tabular}{l|ccccc}
  
  % \cline{1-6}
    \toprule
   \textbf{Experiment} & \textbf{\# Epochs/Steps} & \textbf{Learning Rate} & \textbf{Batch Size} & \textbf{GPUs} & \textbf{Saving Mechanism}\\ \midrule
    \multicolumn{6}{c}{\textbf{Pretraining}}  \\ \midrule
   First-stage/ablations & 300K steps & 1e-4 & 160 & 4xA100 82GB & each 50K steps\\
   Second-stage & 100K steps & 1e-4 & 80 & 4xA100 82GB & each 50K steps\\ \midrule
   
    \multicolumn{6}{c}{\textbf{Finetuning (main 960x960 model)}}  \\\midrule
   ChartQA & 20 epochs & 5e-5 & 24 & 4xV100 32GB & each 1 epoch\\ 
   Chart-to-text Pew & 200 epochs & 5e-5 & 48 & 4xA100 40GB & each 5 epochs\\ 
   Chart-to-text Statista & 100 epochs & 5e-5 & 48 & 4xA100 40GB & each 5 epochs\\ 
   OpenCQA & 200 epochs & 5e-5 & 24 & 4xV100 32GB & each 5 epochs\\
    \bottomrule


 \end{tabular}}
 \label{tab:finetuning-details}
\end{table*}


\subsection{Ablation study}
\label{app:ablation}

To further assess the impact of our different pretraining objectives on our model's performance, we conducted ablation studies. Due to computational limitations, we focused on pretraining the model only on the lower image size (512x512) and compared it against the corresponding main model (512x512). 
From \Cref{tab:ablations}, we observe that removing the Chart Summarization or Open-ended Question Answering objectives  led to a slight decrease in performance on ChartQA. We attributed this to the abundance of numerical reasoning examples in pretraining. However, removing the Numerical Reasoning pretaining task led to a substantial decrease in performance on ChartQA, highlighting the importance of this task in imbuing numerical abilities into our model. Pretraining the model without the Data Table Generation objective resulted in a relatively weak performance in the ChartQA benchmark, underscoring the importance of understanding underlying data tables of charts in answering reasoning questions.

\subsection{Human and ChatGPT Evaluation}
\label{app:human-chatgpt-eval}

\begin{figure*}[t!]
     \centering
        \includegraphics[width=\textwidth]{Figures/ChatGPT_Evaluation.drawio4.pdf}
         \caption{\small{
         The pipeline designed for the ChatGPT Evaluation Experiment. First, we feed the task description followed by our desired criteria into ChatGPT in order to get detailed grading instructions. Then, the chart (underlying data table representation) and a sample summary are appended to the prompt which is fed again into ChatGPT to receive the feedback. 
         }
         }
    \label{fig:chatgpt-eval}
    % }
\end{figure*} 

\begin{figure*}
     \centering
     \begin{subfigure}[b]{\textwidth}
         \centering
         \includegraphics[width=\textwidth]{Figures/Human_eval1.png}
         \caption{}
     \end{subfigure}
     \hfill
     \begin{subfigure}[b]{\textwidth}
         \centering
         \includegraphics[width=\textwidth]{Figures/Human_eval2.png}
         \caption{}
     \end{subfigure}
        \caption{Human evaluation interface.}
        \label{fig:human-eval-interface}
\end{figure*}



\begin{table}[t!]
 \setlength\extrarowheight{2pt}
 \centering
 \caption{
  Human evaluation on summaries for 150 random samples from Chat2text Statista test split.
 }
 \scalebox{0.59}{\begin{tabular}{l|cccc}
  
  \toprule
   \textbf{Criteria} & \textbf{ZeroShot} & \textbf{Finetuned} & \textbf{MatCha} & \textbf{Gold}\\ \midrule
   Factually incorrect sents & 13.45\% & 9.63\% & 21.97\% & \textbf{3.59\%}\\
   \midrule
   Elemental/encoded sents & 19.42\% & 26.06\% & \textbf{29.61\%} & 21.07\% \\ 
   Statistical/relational sents & \textbf{57.41\%} & 33.42\% & 34.07\% & 34.70\% \\ 
   Perceptual/cognitive sents & \textbf{6.98\%} & 1.41\% & 0.31\% & 5.39\% \\
    Contextual/domain-specific sents & 1.36\% & 14.44\% & 7.32\% & \textbf{20.56\%}  \\ \bottomrule
   

 \end{tabular}}
 \label{tab:human_eval}
\end{table}

As discussed in section \Cref{subsec:summary_eval}, we evaluate the following three criteria in the human evaluation study: \emph{(1) Informativeness} which measures how much information from the chart the summary covers. Ideally, an informative summary should contain high-level insights from the chart, such as important patterns, trends, and outliers in data; \emph{(2) Factual Correctness} which considers how accurate the summary is. A factually correct summary should only contain information (e.g. numbers, events, entities) that is true and/or supported by the chart; \emph{(3) Semantic Levels defined by \cite{lundgard2021accessible}} which categorize the content of summaries across four levels: visual encoding (e.g., axis, legends, color), statistical/relational  (e.g., min, max, avg.), perceptual/cognitive (e.g., describing overall trends, complex patterns, or outliers), and context/domain specific information. Our process for evaluating the informativeness is explained in \Cref{subsec:summary_eval}. For factual correctness and semantic level measures, the annotator goes through each sentence of the summary to determine whether the sentence contains any factual error and what levels of semantic content are present for that sentence.
\Cref{tab:human_eval} shows the results of our human evaluation study on factual correctness, and different semantic levels.

\Cref{fig:chatgpt-eval} shows an overview of the paradigm we use in our ChatGPT-driven evaluation study. \Cref{fig:human-eval-interface} depicts the interface we used in our human evaluation study. 


\subsection{Error Analysis}
\label{app:error-analysis}

\Cref{fig:error-analysis} shows the models performance on challenging samples. Q1 and Q2 examples are two visual numerical reasoning questions about charts which look challenging for SoTA models. Q3 is an example of an overpopulated chart with so many data elements which confuses the model to generate insightful summary. Finally, Q4 shows a factual error in a generated summary from finetuned \model.




\begin{table}[t!]
 \setlength\extrarowheight{2pt}
 \centering
 \caption{
  \model\ ablations on ChartQA benchmark.
 }
 \scalebox{0.73}{\begin{tabular}{l|ccc}
  
  % \cline{1-4}
  \toprule
  
  \multirow{2}{*}{} & \multicolumn{3}{c}{ChartQA} \\ \cmidrule(lr){2-4}
  
   Model & aug. & human & avg. \\
                                            
    \midrule
  \model\ (512x512) & 85.84 & 43.60 & 64.72 \\ 
  
  No Chart Summarization & 84.96 & 42.72 & 63.84 \\ 

  No Open-ended Question Answering & 85.52 & 42.96 & 64.24 \\   
  No Numerical \& Visual Reasoning & 84.08 & 35.44 & 59.76 \\ 
  No Data Table Generation & 83.84 & 42.24 & 63.04 \\   
 \bottomrule
 
 \end{tabular}}
 \label{tab:ablations}
\end{table}



\begin{figure*}
    \centering
    \includegraphics[width=0.99\textwidth]{Figures/error_analysis.drawio4.pdf}
    \caption{Some challenging examples from the ChartQA and Chart-to-text benchmarks.}
    \label{fig:error-analysis}
\end{figure*}









\subsection{Time and Memory Efficiency}
\label{app:efficiency}
To compare the time efficiency, we measure the average inference time of the models on three benchmarks: ChartQA, Chart-to-Text (Pew), and Chart-to-Text (Statista) using 10 random samples from each benchmark. The experiments were conducted on Google's Colab platform with cpu type. Overall, \model\ shows much faster inference times compared to MatCha as shown in \Cref{fig:inference_time}. 


\subsection{Templates for Numerical and Visual Reasoning Question Generation}
\label{app:template-visual-logical}
\Cref{tab:templates} is the list of the templates we used to generate numerical and visual reasoning questions.


\begin{table*}[t]
 \centering
 \caption{
  Numerical \& Visual reasoning templates.
 }
 \scalebox{0.54}{\begin{tabular}{|l|}
  
  \toprule
    \textbf{Template-based Questions}\\
    \midrule
   1) First bar from the top/left in the second group from the top/left \\2) First bar from the bottom/right in the second group from the bottom/right \\3) Second bar from the bottom/right in the first group from the bottom/right \\4) Second bar from the right/bottom in the first group from the left/top \\5) Topmost/Leftmost bar \\6) Bottommost/Rightmost bar \\7) Second bar from the top/left \\8) Second bar from the right/bottom \\9) Leftmost topmost bar \\10) Leftmost bottommost bar \\11) Rightmost topmost bar \\12) Rightmost bottommost bar \\13) Leftmost \texttt{<color>} data \\14) Rightmost \texttt{<color>} data \\15) Second from the left \texttt{<color>} data \\16) Second from the right \texttt{<color>} data \\17) Which legend represented by \texttt{<color>}? \\18) What is the color of \texttt{<legend>}? \\19) Which one is greater, \texttt{<x1>} or \texttt{<x2>}? \\20) Divide the sum of largest and lowest values by \texttt{<n>} \\21) When did line \texttt{<legend - label>} peak? \\22) What is the difference between maximum and minimum of \texttt{<legend - label>}? \\23) Sum pie segments above \texttt{<value>} \\24) What is the sum of top three values? \\25) What is the median/mode of \texttt{<legend - label>}? \\26) What is the negative peak of \texttt{<legend - label>}? \\27) What is the largest/smallest value of \texttt{<legend - label>}? \\28) Which two x-axis labels of \texttt{<legend - label>} sums up to \texttt{<value>}? \\29) What is the sum of the second highest and second lowest value of \texttt{<legend - label>}? \\30) Which x-axis label is second highest for \texttt{<legend - label>}? \\31) What is the sum of two middle values of \texttt{<legend - label>}? \\32) Which two x-axis labels of \texttt{<legend - label>} have a difference of \texttt{<value>} ? \\33) What is the average of  \texttt{<legend - label>} from \texttt{<x - label - 1>} to \texttt{<x - label - 2>}? \\34) What is the average of the highest and lowest value of \texttt{<legend - label - l>}? \\35) What is the sum of the average of \texttt{<legend - label - 1>} and average of \texttt{<legend - label - 2>}? \\36) What is the sum/difference of the maximum of \texttt{<legend - label - 1>} and minimum of \texttt{<legend - label - 2>}? \\37) Which x-axis label has the maximum/minimum difference between \texttt{<legend - label - 1>} and minimum of \texttt{<legend - label - 2>}? \\38) Which x-axis label witnessed the smallest value of \texttt{<legend - label>}? \\39) Which label contains largest/smallest values across all labels? \\40) Sum up the medians of all the data series in this chart \\41) What is the average of all values above \texttt{<value>}?  \\42) What is the sum of the largest and smallest difference between \texttt{<legend - label - 1>} and \texttt{<legend - label - 2>}? \\43) What is the maximum/minimum difference between \texttt{<legend - label - 1>} and \texttt{<legend - label - 2>}? \\44) What is the ratio of the largest to the smallest pie segment? \\45) What is the ratio of the two largest/smallest segments? \\46) What is the difference between the leftmost and rightmost bars? \\47) What is the sum of the bars in the second group from the left? \\48) What is the sum of the bars in the first group from the right? \\49) What is the ratio between the two leftmost bars? \\50) What is the difference between the rightmost \texttt{<color - 1>} bar and leftmost \texttt{<color - 2>} bar? \\51) What is the average of  \texttt{<color>} bars values? \\52) How many  \texttt{<color>}bars are larger than \texttt{<N>} ? \\53) What is the average of the bars in the second group from the right? \\54) How many bars in the leftmost group have a value over \texttt{<N>}? \\55) What does the  \texttt{<color>} represent? \\56) What is the median value of the  \texttt{<color>} bars/line? \\57) What is the average of the  \texttt{<color - 1>} sum and  \texttt{<color - 2>} sum? \\58) What is the average of the  \texttt{<color - 1>} median and  \texttt{<color - 2>} median? \\59) What is the least difference between the  \texttt{<color - 1>} and \texttt{<color - 2>} bars/line? \\60) What is the ratio between the leftmost and rightmost bar in the first group from the left? \\61) What is the maximum value in the  \texttt{<color>} bars/line? \\62) What is the minimum value in the  \texttt{<color>} bars/line? \\63) What is the sum of  \texttt{<color>} bars/line? \\64) What is the difference between the maximum values of the two leftmost bar groups? \\65) Sum of the first \texttt{<color - 1>} and last \texttt{<color - 2>} bars/line points \\66) Difference between the two lowest \texttt{<color>} bars \\67) Add largest and smallest \texttt{<color>} line/bar values and divide by 2 \\68) What is the value of \texttt{<color>} line/bars in \texttt{<x - axis - label>}? \\69) Sum/Average of \texttt{<color - 1>} and \texttt{<color - 2>} values in \texttt{<x - axis - label>}? \\70) Sum of highest points in \texttt{<color - 1>} and \texttt{<color - 2>} lines/bars \\71) Which color has the highest/smallest values? \\72) How many values are equal in \texttt{<color - 1>} line/bar? \\73) Sum two rightmost values of \texttt{<color>} graph \\74) Product of two smallest values in the graph \\75) Sum of lowest and median values of \texttt{<color>} graph/bars \\76) When did \texttt{<color>} line reached the peak? \\77) What is the average of the rightmost three points of \texttt{<color>} line? \\78) How many \texttt{<color>} data points are above \texttt{<value>}? \\79) What's the ratio of the largest and the third/second-largest \texttt{<color>} bar? \\80) Is the sum of lowest value of \texttt{<color - 1>} and \texttt{<color - 2>} bar greater than largest value of \texttt{<color - 3>} bar? \\81) Is the median value of \texttt{<color - 1>} bars greater than the median value of \texttt{<color - 2>} bars? \\82) Is the median of all the \texttt{<color - 1>} bars greater than the largest value of \texttt{<color - 2>} bar? \\83) What's the product of \texttt{<color>} bars in India and Japan? \\84) Is the sum of the two middle bars greater than the sum of top and bottom bars? \\85) What's the ratio of the \texttt{<x - axis - label - 1>} \texttt{<color - 1>} bar and the \texttt{<x - axis - 2>} \texttt{<color - 2>} bar? \\86) Is the total of all \texttt{<color - 1>} bars greater than the total of all \texttt{<color - 2>} bars? \\87) Take the sum of the two smallest \texttt{<color - 1>} bars and smallest \texttt{<color - 2>} bars, deduct the smaller value from the larger value, what's the result? \\88) What is the sum/average of two smallest/largest \texttt{<color>} bars? \\89) What is the ratio of \texttt{<color - 1>} and \texttt{<color - 2>} segments? \\90) What segment is represented by \texttt{<color>}? \\
                                             
 \bottomrule
 
 \end{tabular}}
 \label{tab:templates}
\end{table*}